\chapter{Обзор существующих решений по разработке фреймворков для микроконтроллеров}

\label{cha:theory}

Разработка программного обеспечения для микроконтроллеров требует учёта специфики аппаратных ограничений, включая объёмы памяти, энергоэффективность и поддержку периферийных устройств\cite{Harris2021}. Данная глава посвящена теоретическому анализу базовых концепций, принципов и технологий, лежащих в основе разработки фреймворков для микроконтроллеров.
\section{Микроконтроллеры в современных встраиваемых системах}

Микроконтроллеры представляют собой ключевой компонент современных встраиваемых систем, обеспечивая интеллектуальное управление различными устройствами от бытовой техники до промышленных систем управления\cite{Kernighan2019}. Первый коммерчески доступный микроконтроллер TMS1000 объединил в себе 4-битный ЦПУ, 1 КБ ПЗУ, 256-битную ОЗУ и линии ввода/вывода на одном чипе. Сегодня высокопроизводительный 32-битный микроконтроллер с 2 МБ встроенной SRAM может достигать скорости 1 ГГц, работая как минимум в 2500 раз быстрее. Современные микроконтроллеры часто классифицируют по разрядности обрабатываемых данных: 8-разрядные, 16-разрядные и 32-разрядные.

В зависимости от функционального назначения микроконтроллеры могут быть разделены на следующие основные типы:

\begin{itemize}
	\item Встраиваемые 8-разрядные микроконтроллеры
	\item 16- и 32-разрядные микроконтроллеры
	\item Цифровые сигнальные процессоры (DSP)
\end{itemize}

Существуют две основные архитектуры микроконтроллеров: архитектура фон Неймана и гарвардская архитектура\cite{MartinArch2024}. Гарвардская архитектура, характеризующаяся разделением каналов для инструкций и данных, обеспечивает более высокую производительность за счет одновременного доступа к командам и данным (\ref{fig:arch_comp}).

\begin{figure}
	\centering
	\includegraphics[width=0.7\textwidth]{inc/img/arch-comp}
	\caption{Сравнение архитектур}
	\label{fig:arch_comp}
\end{figure}

В гарвардской архитектуре характеристики устройств памяти для инструкций и данных могут различаться. В частности, ширина слова, тактирование, технология реализации и структура адресов памяти не обязательно должны быть одинаковыми. Это позволяет оптимизировать каждый канал под конкретные требования, что особенно важно для встраиваемых систем с ограниченными ресурсами.

Современные микроконтроллеры часто используют гибридные архитектуры, сочетающие преимущества обоих подходов. Например, многие CISC-процессоры имеют раздельную кэш-память первого уровня для команд и данных, что позволяет им получать одновременно и команду, и данные для её выполнения.

Встраиваемые микроконтроллеры содержат значительное число вспомогательных устройств, что позволяет использовать минимальное количество дополнительных компонентов при интеграции в систему. Типичная структура микроконтроллера представлена на рисунке \ref{fig:mik_struct}.

\begin{figure}
	\centering
	\includegraphics[width=0.7\textwidth]{inc/img/mik-struct}
	\caption{Структурная схема микроконтроллера}
	\label{fig:mik_struct}
\end{figure}

В состав современного микроконтроллера обычно входят следующие компоненты:

\begin{itemize}
	\item Центральный процессор (CPU)
	\item Память программ (Flash, ROM, EPROM)
	\item Память данных (RAM)
	\item Схема начального запуска (Reset)
	\item Генератор тактовых импульсов
	\item Периферийные устройства (таймеры, АЦП, ЦАП, интерфейсы связи)
\end{itemize}

\section{Архитектура RISC-V и её применение в микроконтроллерах}

RISC-V представляет собой открытую стандартную архитектуру набора инструкций, основанную на принципах RISC\cite{RISC2024}. В отличие от проприетарных архитектур, таких как ARM, AVR или x86, RISC-V является открытым стандартом, что позволяет любому производителю реализовывать его спецификации без лицензионных отчислений.

\Abbrev{RISC}{reduced instruction set computer ""--- архитектурный подход к проектированию процессоров, в которой быстродействие увеличивается за счёт такого кодирования инструкций, чтобы их декодирование было более простым, а время выполнения -- меньшим.}

RISC-V был разработан в Калифорнийском университете в Беркли, начиная с 2010 года, как перспективная архитектура, не обремененная требованиями обратной совместимости\cite{RISC2024}. Она была спроектирована не только для академических целей, но и для практического применения в реальных устройствах.

Преимущества архитектуры RISC-V для применения в микроконтроллерах:

\begin{itemize}
	\item Модульность: RISC-V состоит из базового набора инструкций (RV32I, RV64I) и опциональных расширений (M, A, F, D, C и др.), что позволяет адаптировать процессор под конкретные задачи\cite{RISC2024}.
	\item Открытость: Спецификации RISC-V доступны бесплатно под лицензией BSD, что способствует инновациям и снижает барьеры для входа новых производителей.
	\item Масштабируемость: RISC-V поддерживает 32-, 64- и 128-битные варианты адресного пространства, что обеспечивает гибкость при реализации различных классов устройств.
	\item Компактность: Расширение "C" позволяет использовать 16-битные инструкции, сокращая размер кода до 25-30\% по сравнению с базовым набором.
	\item Архитектура load-store: Инструкции RISC-V адресуют только регистры, а для взаимодействия с памятью используются специальные инструкции загрузки и сохранения.
\end{itemize}

RISC-V предлагает гибкую систему расширений набора инструкций, что позволяет создавать процессоры с различными характеристиками. Основные наборы инструкций включают:

\begin{itemize}
	\item RV32I - базовый 32-битный набор целочисленных инструкций
	\item RV64I - базовый 64-битный набор целочисленных инструкций
	\item RV32E - набор с уменьшенным количеством регистров для встраиваемых систем
\end{itemize}

Стандартные расширения включают:

\begin{itemize}
	\item I (Base Integer Instruction Set): Базовый набор целочисленных инструкций, обязательный для всех реализаций RISC-V.
	\item M (Standard Extension for Integer Multiplication and Division): Стандартное расширение для операций умножения и деления целых чисел.
	\item A (Standard Extension for Atomic Instructions): Стандартное расширение для атомарных операций, необходимых для синхронизации в многопоточных приложениях.
	\item F (Standard Extension for Single-Precision Floating-Point): Стандартное расширение для операций с плавающей точкой одинарной точности (32 бита).
	\item D (Standard Extension for Double-Precision Floating-Point): Стандартное расширение для операций с плавающей точкой двойной точности (64 бита).
	\item C (Standard Extension for Compressed Instructions): Стандартное расширение для 16-битных сжатых инструкций, позволяющее уменьшить размер кода и повысить плотность инструкций.
	\item V (Vector Extension): Стандартное расширение для векторных вычислений, предназначенное для ускорения обработки данных в приложениях, таких как машинное обучение и обработка сигналов.
\end{itemize}

Полная спецификация ядра обозначается сочетанием базового набора и расширений, например, RV32IMC означает 32-битное ядро с поддержкой целочисленного умножения/деления и сжатых инструкций. На рисунке \ref{fig:rv32imac} представлен основной набор команд.

\begin{figure}
	\centering
	\includegraphics[width=0.7\textwidth]{inc/img/rv32imac}
	\caption{Набор команд RISC-V}
	\label{fig:rv32imac}
\end{figure}

Применение RISC-V в микроконтроллерах становится всё более распространённым благодаря преимуществам открытой архитектуры. Компании, такие как Renesas, SiFive и GigaDevice, активно разрабатывают микроконтроллеры на базе RISC-V для различных приложений\cite{RISC2024}. Примером микроконтроллерного ядра на базе RISC-V является SCR1 от компании Syntacore. Это открытый и проверенный в кремнии RISC-V-совместимый 32-битный микроконтроллерный ядро начального уровня, ориентированный на встраиваемые приложения. Основные характеристики SCR1 включают:

\begin{itemize}
	\item Поддержка RISC-V стандарта I, E, M и C расширений
	\item 4-стадийный конвейер
	\item Блок TCM (Tightly-Coupled Memory) размером до 64KB
	\item Блок IPIC для обработки прерываний
	\item Стандартные интерфейсы AXI4, AHB-Lite и JTAG/cJTAG
	\item Два счётчика производительности и встроенный 64-битный RTC таймер
	\item Ядро SCR1 используется в отечественном микроконтроллере MIK32 АМУР (К1948ВК018), который является первым российским микроконтроллером с архитектурой RISC-V.
\end{itemize}

Архитектура RISC-V предоставляет ряд преимуществ для микроконтроллеров:

\begin{itemize}
	\item Экономическая составляющая: Отсутствие лицензионных отчислений снижает стоимость разработки и производства.
	\item Гибкость проектирования: Модульная структура позволяет создавать оптимизированные решения для конкретных приложений.
	\item Энергоэффективность: Простота и оптимизированный набор инструкций способствуют снижению энергопотребления.
	\item Безопасность: Открытость архитектуры позволяет проводить полный аудит безопасности реализации.
\end{itemize}

В таблице \ref{tab:microcontroller_comparison} приведено сравнение двух микроконтроллеров с разными архитектурами.

\begin{center}
	\begin{longtable}{|p{0.28\textwidth}|p{0.32\textwidth}|p{0.32\textwidth}|}
		\caption{Сравнение микроконтроллеров MIK32 АМУР и ARM Cortex-M0+ (STM32G030F6P6)}
		\label{tab:microcontroller_comparison}
		\\ \hline
		\textbf{Характеристика} & \textbf{MIK32 АМУР} & \textbf{ARM Cortex-M0+ (STM32G030F6P6)} \\
		\hline \endfirsthead
		\subcaption{Продолжение таблицы~\ref{tab:microcontroller_comparison}} \\
		\hline
		\textbf{Характеристика} & \textbf{MIK32 АМУР} & \textbf{ARM Cortex-M0+ (STM32G030F6P6)} \\
		\hline \endhead
		\subcaption{Продолжение на след. стр.} \endfoot
		\hline \endlastfoot
		Архитектура & RISC-V & ARMv6-M \\
		\hline
		Максимальная тактовая частота & 32 МГц & 64 МГц \\
		\hline
		ОЗУ & 16 Кбайт & 8 Кбайт \\
		\hline
		ПЗУ & 8 Кбайт EEPROM, 256/512 бит OTP & 32 Кбайт Flash \\
		\hline
		Внешняя память & До 16 МБ через QSPI Flash & Зависит от реализации \\
		\hline
		SPI интерфейс & 2 модуля & 1 модуль \\
		\hline
		I2C интерфейс & 2 контроллера & 1 контроллер \\
		\hline
		UART интерфейс & 3 порта & 2 порта \\
		\hline
		USB интерфейс & Отсутствует & FS 2.0 (зависит от реализации) \\
		\hline
		Аппаратное шифрование & ГОСТ 34.12-2018, AES128 & Отсутствует \\
		\hline
		Генератор случайных чисел & Да & Нет \\
		\hline
		Таймеры с ШИМ & 3×32 бит, 3×16 бит с ШИМ & Зависит от реализации (количество не указано) \\
		\hline
	\end{longtable}
\end{center}

\section{Обзор существующих фреймворков и RTOS для микроконтроллеров}

Фреймворк для микроконтроллера представляет собой программную платформу, обеспечивающую абстракцию над аппаратным уровнем и предоставляющую набор инструментов и API для разработки приложений\cite{MartinCleanCode2024}. Операционная система реального времени (RTOS) расширяет возможности фреймворка, добавляя механизмы планирования задач и управления ресурсами с гарантированным временем отклика.

Основные компоненты современного фреймворка для микроконтроллеров включают:

\begin{itemize}
	\item Слой абстракции оборудования (HAL)
	\item Драйверы периферийных устройств
	\item Библиотеки для работы с коммуникационными протоколами
	\item Средства управления питанием
	\item Инструменты разработки и отладки
\end{itemize}

В случае RTOS добавляются также:
\begin{itemize}
	\item Планировщик задач
	\item Средства синхронизации (семафоры, мьютексы, очереди)
	\item Система управления памятью
	\item Обработчики прерываний
\end{itemize}

\subsection{FreeRTOS}

FreeRTOS является одной из наиболее популярных операционных систем реального времени для микроконтроллеров\cite{Harris2021}. Она имеет открытый исходный код и распространяется под лицензией MIT с 2017 года.

FreeRTOS занимает лидирующие позиции в сегменте операционных систем реального времени для ресурсоограниченных устройств. Её архитектура ориентирована на минимальное потребление ресурсов: для работы требуется менее 10 КБ оперативной памяти и 20 КБ постоянной памяти. Основная функциональность включает планировщик задач, механизмы синхронизации (семафоры, очереди, мьютексы) и поддержку прерываний.

Особенностью FreeRTOS является интеграция с платформами, такими как ESP-IDF для микроконтроллеров ESP32. Например, в ESP32 задачи FreeRTOS используются для параллельного выполнения сетевых операций и управления периферией. Принцип работы планировщика показан на рисунке \ref{fig:TaskExecution}.

\begin{figure}
	\centering
	\includegraphics[width=0.7\textwidth]{inc/img/TaskExecution}
	\caption{Схема выполнения задач во FreeRTOS}
	\label{fig:TaskExecution}
\end{figure}

FreeRTOS портирована на более чем 35 микропроцессорных архитектур, включая ARM (ARM7, ARM9, Cortex-M0, Cortex-M3, Cortex-M4, Cortex-A), Atmel AVR, AVR32, PIC, MSP430 и RISC-V\cite{Harris2021}. Это делает её универсальным выбором для широкого спектра микроконтроллеров.

\subsection{Zephyr}

Zephyr Project представляет собой легкую, масштабируемую операционную систему реального времени, разработанную консорциумом компаний (Intel, NXP, Synopsys, UbiquiOS) под эгидой Linux Foundation\cite{MartinArch2024}. Она распространяется под лицензией Apache 2.0 и предназначена для устройств с ограниченными ресурсами.

Zephyr отличается модульной архитектурой, которая позволяет разработчикам включать только необходимые компоненты:

\begin{itemize}
	\item Микроядро: Обеспечивает базовые функции, такие как управление памятью, задачами и обработку прерываний в реальном времени
	\item Модульность: Возможность добавления или исключения компонентов делает Zephyr гибким для различных приложений
	\item Широкая поддержка архитектур: Совместимость с ARM, x86, RISC-V и другими платформами упрощает переносимость приложений
\end{itemize}

Основные особенности Zephyr включают:
\begin{itemize}
	\item Единое адресное пространство (монолитный образ)
	\item Определение ресурсов во время компиляции
	\item Минимальный контроль ошибок для оптимизации производительности
\end{itemize}

\subsection{Mbed OS}

Mbed OS представляет собой открытую операционную систему для Интернета вещей (IoT) на базе микроконтроллеров Cortex-M\cite{MartinCleanCode2024}. Она разработана компанией Arm и предоставляет абстракцию для микроконтроллеров, позволяя писать приложения на C/C++, которые работают на любой Mbed-совместимой плате. Mbed OS ориентирована на IoT-приложения и включает в себя следующие компоненты:

\begin{itemize}
	\item Аппаратно-независимый слой (HAL): Упрощает взаимодействие с периферией через стандартные API (SPI, I2C, UART)
	\item Поддержка хранения данных: Включает файловые системы (FAT, NOR flash)
	\item Интеграция сетевых технологий: Поддержка Bluetooth LE, Wi-Fi, LoRaWAN и других протоколов
	\item Полный профиль: Включает RTOS (Keil RTX) и все API
	\item Bare metal профиль: Без RTOS, для приложений, не требующих сложного управления потоками
\end{itemize}

Mbed OS также предлагает встроенные средства безопасности (Mbed TLS) для защиты данных и обеспечения безопасных коммуникаций в IoT-устройствах.

\section{Сравнение фреймворков для встраиваемых систем}

Результаты сравнительного анализа показывают, что FreeRTOS является наиболее подходящим кандидатом для адаптации к MIK32 благодаря минимальным требованиям к ресурсам и хорошей поддержке RISC-V\cite{Harris2021}. Однако для полноценного использования потребуется разработка специфичных для MIK32 драйверов периферийных устройств и оптимизация системы обработки прерываний.

Zephyr также представляет интерес благодаря модульной структуре и поддержке RISC-V, но его требования к памяти могут быть критичны для MIK32\cite{MartinArch2024}. Mbed OS, будучи ориентированным на микроконтроллеры Arm Cortex-M, потребует значительных усилий для адаптации к архитектуре RISC-V. На основе проведенного анализа составлена сравнительная таблица \ref{tab:frameworks_comparison} основных характеристик рассмотренных фреймворков и RTOS:

\begin{center}
	\begin{longtable}{|p{0.22\textwidth}|p{0.22\textwidth}|p{0.22\textwidth}|p{0.22\textwidth}|}
		\caption{Сравнение фреймворков для микроконтроллеров}
		\label{tab:frameworks_comparison}
		\\ \hline
		\textbf{Критерий} & \textbf{FreeRTOS} & \textbf{Zephyr} & \textbf{Mbed OS} \\
		\hline \endfirsthead
		\subcaption{Продолжение таблицы~\ref{tab:frameworks_comparison}} \\
		\hline
		\textbf{Критерий} & \textbf{FreeRTOS} & \textbf{Zephyr} & \textbf{Mbed OS} \\
		\hline \endhead
		\subcaption{Продолжение на след. стр.} \endfoot
		\hline \endlastfoot
		Потребление энергии & Управление & Оптимизация через доступ к железу & Энергосберегающие API \\
		\hline
		Поддержка периферии & Требует внешних библиотек & Широкий набор встроенных драйверов & Обширная библиотека драйверов \\
		\hline
		Сетевые протоколы & Через дополнительные библиотеки & Встроенная поддержка & Встроенная поддержка \\
		\hline
		Размер кода & Минимальный & Модульный, зависит от конфигурации & Больше, чем у FreeRTOS \\
		\hline
		Лицензия & MIT & Apache 2.0 & Apache 2.0 \\
		\hline
		Поддержка RISC-V & Да & Да & Частичная \\
		\hline
		Экосистема & Обширная & Растущая & Развитая для Arm \\
		\hline
	\end{longtable}
\end{center}

Проведенный анализ существующих решений для разработки фреймворков микроконтроллеров позволяет сделать следующие выводы:

\begin{itemize}
	\item Архитектура RISC-V предоставляет значительные преимущества для создания отечественных микроконтроллеров благодаря открытости спецификации, модульности набора инструкций и отсутствию лицензионных ограничений. Микроконтроллер MIK32 на базе ядра SCR1 демонстрирует конкурентоспособные характеристики по сравнению с зарубежными аналогами.
	\item Анализ существующих фреймворков показал, что FreeRTOS является наиболее подходящим базисом для адаптации к MIK32 благодаря минимальным требованиям к ресурсам и хорошей поддержке RISC-V. Zephyr представляет интерес благодаря модульной архитектуре, однако требования к памяти могут быть критичны. Mbed OS требует значительных усилий для адаптации к архитектуре RISC-V.
	\item Отсутствие специализированного фреймворка для MIK32 создает необходимость разработки оригинального решения, учитывающего особенности архитектуры микроконтроллера и требования отечественного рынка встроенных систем. Ключевыми требованиями к такому фреймворку являются минимальное потребление ресурсов, модульная архитектура и поддержка специфических периферийных устройств MIK32.
\end{itemize}


\Abbrev{MCU}{microcontroller unit ""--- микроконтроллер, интегральная схема, предназначенная для управления электронными устройствами, объединяющая процессор, память и периферию на одном чипе}

\Abbrev{HAL}{hardware abstraction layer ""--- слой абстракции оборудования, предоставляющий унифицированный интерфейс для взаимодействия с периферийными устройствами микроконтроллера}

\Abbrev{RTOS}{real-time operating system ""--- операционная система реального времени, обеспечивающая гарантированное время отклика на внешние события}

\Abbrev{CLI}{command line interface ""--- интерфейс командной строки, позволяющий взаимодействовать с системой через текстовые команды}

\Abbrev{SDK}{software development kit ""--- комплект программных инструментов для разработки приложений под определённую платформу}

\Abbrev{UART}{universal asynchronous receiver-transmitter ""--- универсальный асинхронный приёмопередатчик, интерфейс для последовательной передачи данных}

\Define{Микроконтроллер}{устройство на интегральной микросхеме, включающее центральный процессор, память и периферийные интерфейсы, предназначенное для управления электронными системами}

\Define{Слой абстракции оборудования}{программный слой, предоставляющий стандартные интерфейсы для работы с периферией микроконтроллера, скрывая детали реализации аппаратной части}

\Define{Планировщик задач}{компонент RTOS, отвечающий за распределение времени процессора между задачами в соответствии с выбранной стратегией планирования}

\Define{Драйвер периферийного устройства}{программный модуль, реализующий взаимодействие с конкретным аппаратным компонентом микроконтроллера (например, UART, SPI, I2C)}

\Abbrev{EEPROM}{electrically erasable programmable read-only memory ""--- энергонезависимая память, допускающая многократную перезапись данных}
